\section{Introduction}
\label{sec:intro}

A latent quantity $\Y$ is often observed only through a coarsening: a mechanism
that maps $\Y$ to a coarser report $R$ and discards the rest, so that the report
identifies only a \emph{candidate set} $c(R)$ of latent values consistent with it.
A failed component is reported as a set of suspects; a transcript count is reported
after a dropout; a sensitive statistic is reported after calibrated noise; a label
is reported as a vote of imperfect heuristics; a disease state is reported as a
billing code. The coarsening-at-random conditions C1, C2, C3
\citep{heitjan1991ignorability,gill1997coarsening,little2002statistical} say when
the analyst may ignore the coarsening mechanism and maximize a face-value
likelihood, and a companion synthesis \citep{towell2026synthesis} shows that, when
they hold, one consistency identity and one rank condition govern identifiability
across six fields, with a \emph{singleton} candidate set ($|c(R)| = 1$) restoring
identifiability whenever it fails.

That program assumes its conditions. The most fragile of the three is C2, the
requirement that within a candidate set the coarsening probability does not depend
on which admissible value is the true one. C2 is precisely the
coarsened-data analogue of missing at random, and like missing at random it is
\emph{informative when it fails and untestable from the observed data}
\citep{grunwald2003updating,jaeger2005ignorability}. It also fails in exactly the
way each application cares about: dropout depends on true expression, coding depends
on true severity, data-dependent privacy mechanisms depend on the realized query,
labeling functions are correlated through the true label, and masking is
group-structured. Treating C2 as benign is therefore the program's main exposure,
and the realistic regime is not C2 but a bounded violation of it.

This paper develops the theory for that regime. We ask two questions. First, when
C2 fails by a controlled amount, how wrong is the face-value fit, and what is the
identified set for the latent parameter? Second, how much external information, in
the form of singletons, is needed to buy identification back? The answers place the
coarsening program in two older and deeper lineages it has not yet engaged:
missing-not-at-random sensitivity analysis
\citep{copas1997inference,manski1989anatomy,molenberghs2008every,daniels2008missing},
where a bounded departure from ignorability yields an identified set rather than a
point, and measurement-error double sampling and verification-bias correction
\citep{tenenbein1970double,begg1983assessment,carroll2006measurement}, where a
validation subsample, our singleton, restores identification. The singleton is
classical internal validation; the C2 tilt is a structured sensitivity parameter.

\subsection{Contributions}

\begin{enumerate}[(i)]
\item \textbf{A tilt parametrization of C2 violation} (\cref{sec:framework}).
  We relax C2 to a log-linear tilt of magnitude $\delta$ over each candidate set,
  recovering C2 at $\delta = 0$ and the informative-masking model of
  \citet{towell2026mdrelax} as the reliability instance.
\item \textbf{A sensitivity bound} (\cref{thm:sensitivity}, \cref{sec:sensitivity}).
  The asymptotic bias of the face-value MLE is, to leading order, linear in
  $\delta$ with an explicit constant: the inverse face-value information times the
  covariance of the face-value score with the tilt. The latent parameter is
  partially identified, over a set of size order $\delta$ that contracts to a point
  as $\delta \to 0$. In a regular exponential family, when the tilt is linear in the
  natural statistic, the leading term is exact at every $\delta$ along the curved
  directions (the range of $\partial \eta / \partial \theta$, the model's tangent
  space); the marginal-fit check that \citet{towell2026synthesis} shows is blind to
  bias under C2 is exactly the quantity the bound controls.
\item \textbf{A singleton sample-complexity theorem}
  (\cref{thm:restoration}, \cref{sec:restoration}). When the augmented
  candidate-set matrix has rank deficit $r$, a set of singletons spanning the
  deficient subspace restores point identification, and the number needed to
  recover those $r$ directions to a target accuracy is of order $r / \gamma^2$,
  where $\gamma$ is the domain's identification margin. This unifies the gold-set
  rate of \citet{towell2026weaksupcoarsening} and the spike-in-fit bias of
  \citet{towell2026scrnacoarsening} as one rate with a domain-specific margin.
\item \textbf{Six domain instances} (\cref{sec:instances}) and a cross-domain
  simulation (\cref{sec:validation}) showing the bias bound is tight and the
  $r/\gamma^2$ rate holds on representative exponential-family and location-family
  instances.
\end{enumerate}

Together with \citet{towell2026synthesis}, which handles the C2-holds case, this is
the second of two pillars: the synthesis says what is recoverable when coarsening is
ignorable; this paper says how recovery degrades when it is not, and what restores
it.
